% ===== Document Preamble =====
\documentclass[11pt]{article}
\usepackage[margin=0.5in]{geometry}
\usepackage{amsmath,amssymb}
\usepackage{xcolor}
\usepackage{enumitem}
\usepackage{needspace}
\usepackage{comment}

% ===== Colors =====
\definecolor{AnsRed}{RGB}{170,0,0}
\definecolor{WrongBlue}{RGB}{0,0,160}
\definecolor{SkipGray}{RGB}{120,120,120}

% ===== Choice list styling (tight, consistent) =====
\setlist[enumerate]{%
  leftmargin=1.5em,
  itemsep=0.25em,
  topsep=0.25em,
  label=\textbf{\Alph*.}
}

% ===== Style Wrappers =====
\newcommand{\QuestionStyle}[1]{{\large\bfseries\boldmath #1}}
\newcommand{\AnswerStyle}[1]{{\large\bfseries\boldmath #1}}
\newcommand{\ReasonStyle}[1]{{\normalsize #1}}

% ===== Show / Hide Reasoning Toggle =====
\newif\ifshowreason
\showreasontrue      % Instructor version
% \showreasonfalse   % Student version

% ===== Choice Item Macros (Answer + Reasoning) =====
% Usage:
%   \Correct{Answer}{Reasoning}
%   \Wrong{Answer}{Reasoning}

\newcommand{\Correct}[2]{%
  \item \AnswerStyle{\textcolor{AnsRed}{#1}}%
  \ifshowreason\\[-0.2em]
    \ReasonStyle{\textcolor{AnsRed}{#2}}%
  \fi
}

\newcommand{\Wrong}[2]{%
  \item \AnswerStyle{\textcolor{WrongBlue}{#1}}%
  \ifshowreason\\[-0.2em]
    \ReasonStyle{\textcolor{WrongBlue}{#2}}%
  \fi
}

% ===== Multiple-Choice Question Block Macro =====
% Usage: \QBlock{Q1}{stem...}{<items>}{answer key}

\newcommand{\QBlock}[4]{%
  \Needspace{8\baselineskip}%
  \noindent
  \begin{minipage}{\linewidth}
    \QuestionStyle{#1.~#2}

    \vspace{0.25em}
    \begin{enumerate}
      #3
    \end{enumerate}

    \vspace{0.1em}
    \noindent\textcolor{AnsRed}{\textit{\textbf{Answer:}~#4}}

    \vspace{0.25em}
    \hrule height 0.4pt
  \end{minipage}
  \vspace{0.75em}
}

% ===== Open-Ended Question Block Macro =====
% Usage: \OEBlock{Q1}{stem...}{answer text}

\newcommand{\OEBlock}[3]{%
  \Needspace{5\baselineskip}%
  \noindent
  \begin{minipage}{\linewidth}
    \QuestionStyle{#1.~#2}

    \vspace{0.25em}
    \noindent\textcolor{AnsRed}{\textit{\textbf{Answer:}~#3}}

    \vspace{0.25em}
    \hrule height 0.4pt
  \end{minipage}
  \vspace{0.75em}
}

% ===== Multi-Part Question Block Macro =====
% Usage: \MPBlock{Q1}{stem...}{<parts>}

\newcommand{\MPBlock}[3]{%
  \Needspace{8\baselineskip}%
  \noindent
  \begin{minipage}{\linewidth}
    \QuestionStyle{#1.~#2}

    \vspace{0.25em}
    #3

    \vspace{0.25em}
    \hrule height 0.4pt
  \end{minipage}
  \vspace{0.75em}
}

% ===== MC Sub-Part Macro =====
% Usage: \MCPart{a}{stem...}{<items>}{answer key}

\newcommand{\MCPart}[4]{%
  \par\vspace{0.35em}
  \noindent\textbf{(#1)}~#2
  \begin{enumerate}
    #3
  \end{enumerate}
  \vspace{0.1em}
  \noindent\textit{\textbf{Answer:}~#4}
  \par
}

% ===== Open-Ended Sub-Part Macro =====
% Usage: \OEPart{a}{stem...}{answer text}

\newcommand{\OEPart}[3]{%
  \par\vspace{0.35em}
  \noindent\textbf{(#1)}~#2
  \par\vspace{0.15em}
  \noindent\textit{\textbf{Answer:}~#3}
  \par
}

% ===== Skipped Question Note Macro =====
% Usage: \SkipBlock{Q1}{reason}

\newcommand{\SkipBlock}[2]{%
  \noindent\textcolor{SkipGray}{\textit{[#1 omitted: #2]}}%
  \vspace{0.5em}\par
}

% ===== Document Start =====

\begin{document}

\begin{center}
{\huge \textbf{Physics Isomorphic Questions}}\\[0.25em]
{\Large Vectors and Kinematics}\\[0.25em]
{\large Questions generated below. Multiple-choice answers in RED; incorrect answers in BLUE.}
\end{center}

\vspace{0.75em}

\QBlock{Q1}{A survey robot starts at the origin of a coordinate system and detects a beacon at position $\vec r = 8\hat i - 15\hat j$, where distances are measured in meters. Which of the following statements is correct?}{
\Wrong{The distance from the robot to the beacon is $23\,\text{m}$.}{This adds the component magnitudes, $8+15$, instead of finding the vector magnitude. The distance must be computed with the Pythagorean relation, not by simple component addition.}
\Correct{The distance from the robot to the beacon is $17\,\text{m}$.}{The distance from the origin to a point is the magnitude of its position vector. Using $|\vec r|=\sqrt{x^2+y^2}$, we get $|\vec r|=\sqrt{8^2+(-15)^2}=\sqrt{64+225}=\sqrt{289}=17\,\text{m}$. Therefore this statement is correct.}
\Wrong{The distance from the robot to the beacon is $15\,\text{m}$.}{This uses only the $y$-component magnitude and ignores the $x$-component. That is a component-versus-magnitude mistake.}
\Wrong{The distance from the robot to the beacon is $8\,\text{m}$.}{This uses only the $x$-component and ignores the $y$-component. A full 2D distance requires both perpendicular components.}
}{B}

\QBlock{Q2}{The displacement of a particle moving along the $x$-axis is given by $\vec s(t)=\big((2.4\,\text{m/s}^3)t^3-(1.8\,\text{m/s}^2)t^2\big)\hat i$. Which of the following statements is correct?}{
\Wrong{This motion has constant acceleration.}{This comes from not differentiating the cubic displacement term completely. Because $s_x$ contains a $t^3$ term, the acceleration depends on time and is not constant.}
\Wrong{The acceleration at $t=0$ is $-1.8\,\text{m/s}^2$.}{This mistakes the coefficient of the $t^2$ term in displacement for the acceleration itself. You must take two time derivatives before evaluating at $t=0$.}
\Correct{The magnitude of the acceleration at $t=0$ is $3.6\,\text{m/s}^2$.}{Acceleration is the second derivative of displacement. From $s_x=(2.4)t^3-(1.8)t^2$, we get $v_x=\frac{ds_x}{dt}=7.2t^2-3.6t$ and $a_x=\frac{dv_x}{dt}=14.4t-3.6$, so at $t=0$, $a_x=-3.6\,\text{m/s}^2$ and its magnitude is $3.6\,\text{m/s}^2$.}
\Wrong{The magnitude of the acceleration at $t=0$ is $0\,\text{m/s}^2$.}{This usually comes from confusing initial velocity or displacement with acceleration. Evaluating the actual acceleration function at $t=0$ does not give zero.}
}{C}

\QBlock{Q3}{A glider moves from position $\vec r_A = 1\hat i - 7\hat j$ to position $\vec r_B = -5\hat i + 1\hat j$ in $4\,\text{s}$, where positions are measured in meters. Which of the following statements is correct?}{
\Wrong{The distance traveled by the glider is $10\,\text{m}$.}{This confuses the magnitude of displacement with the actual path length traveled. From initial and final positions alone, the total distance traveled cannot be determined.}
\Wrong{The speed at $t=0$ must be $2.5\,\text{m/s}$.}{This incorrectly treats displacement magnitude divided by time as the instantaneous speed at the start. The given data are only enough to find an average based on displacement, not the speed at $t=0$.}
\Wrong{The acceleration at $t=0$ must be $0.625\,\text{m/s}^2$.}{This invents an instantaneous acceleration from only two positions and one time interval. That information is insufficient to determine acceleration at a specific instant.}
\Correct{None of the above is correct.}{The displacement is $\Delta\vec r=\vec r_B-\vec r_A=(-6)\hat i+8\hat j$, so its magnitude is $|\Delta\vec r|=\sqrt{(-6)^2+8^2}=\sqrt{100}=10\,\text{m}$. However, that value is the displacement magnitude, not necessarily the distance traveled, and neither the instantaneous speed at $t=0$ nor the instantaneous acceleration at $t=0$ can be deduced from the given information alone. Therefore none of A, B, or C must be true.}
}{D}

\end{document}
